\documentclass[11pt]{amsart}
\usepackage{lmodern}
\usepackage{amsmath,amssymb,amsthm,booktabs}
\usepackage[hidelinks]{hyperref}

\newtheorem{theorem}{Theorem}
\newtheorem{lemma}[theorem]{Lemma}
\newtheorem{proposition}[theorem]{Proposition}
\newcommand{\Z}{\mathbb Z}
\newcommand{\col}{\operatorname{col}}

\title[Minimal orientations in a seven-cycle family]
 {Minimal orientations in a seven-cycle family of critical graphs}
\author{Alex Chan}
\address{Harvard University and National Bureau of Economic Research}
\email{achan@hbs.edu}
\subjclass[2020]{05C15, 05C20}
\keywords{critical graph, minimal orientation, directed path,
 graph coloring, satisfiability certificate, Hanson--Toft family}
\date{}

\begin{document}

\begin{abstract}
Hanson and Toft observed that a seven-cycle construction gives a
$k$-critical graph for each $k\geq4$.  The first two graphs were known
to have no orientation with exactly one directed path of length $k-1$;
Jensen and Toft asked whether this holds throughout the family.  We
prove that it does.  A short identity for color classes on the cycle
reduces the infinite statement to the six cases $4\leq k\leq9$.  For
these cases we give Boolean encodings of a unique full increasing path,
unsatisfiability certificates, and a small independent checker that
reconstructs each formula and verifies its proof.  The six formula
refutations are also verified in Lean using native evaluation.
\end{abstract}
\maketitle

\section{Introduction}

An orientation of a $k$-chromatic graph is \emph{minimal} if it has
exactly one directed simple path of length $k-1$.  Gallai asked whether
every critical graph has a minimal orientation \cite{Gallai}.  Youngs answered this
question negatively, beginning with a $14$-vertex $4$-critical graph
\cite{Youngs}.  Hanson and Toft identified this graph as the first of a
seven-cycle family $H_k$.  Jensen and Toft recorded the negative cases
$k=4,5$ and asked whether the entire family has no minimal orientation
\cite[\S13.8, pp.~215--216]{JT}.  We answer that family question.

Our proof has two parts.  An elementary identity concerning the seven
attachment triples permits a common color to be deleted from every
clique whenever $k\geq10$.  Every full increasing path of the smaller
graph then extends injectively to one of the larger graph.  The cases
$k=4,\ldots,9$ are finite: their exact unsatisfiable formulas and
checkable proof traces are included in the accompanying supplementary
archive.  The reduction, formula semantics, and certificate format are
specified below so that the finite computation is a verifiable premise
of the theorem.

Write $C=x_0x_1\cdots x_6x_0$, with indices in $\Z/7\Z$.  For each
$i$, take a clique $Q_i$ of order $k-3$, with the seven cliques mutually
disjoint and disjoint from $C$.  Join every vertex of $Q_i$ to
\begin{equation}\label{eq:triples}
 T_i=\{x_i,x_{i+1},x_{i+3}\},
\end{equation}
and add no other edges.  Denote the resulting graph by $H_k$.  Direct
counting gives
\begin{equation}\label{eq:parameters}
 |V(H_k)|=7(k-2),\qquad
 |E(H_k)|=\frac{7}{2}(k^2-k-4),\qquad
 \Delta(H_k)=3k-7.
\end{equation}

An \emph{ordered $k$-coloring} is a proper vertex coloring by the
ordered palette $0,\ldots,k-1$.  A \emph{full increasing path} is a
simple path whose vertices have colors $0,1,\ldots,k-1$ in that order.
We use the following orientation--coloring criterion stated in
\cite[\S13.8]{JT}: a $k$-critical graph admits a minimal orientation if
and only if it has an ordered $k$-coloring with exactly one full
increasing path.  The forward implication applies to \emph{arbitrary}
orientations: Gallai's reduction first gives an acyclic minimal
orientation and then its longest-path levels give the coloring.
Conversely, orienting each edge from its smaller to its larger color
gives the desired acyclic orientation.  Thus it suffices to prove that
every ordered $k$-coloring of $H_k$ has at least two full increasing
paths.

\section{The cycle identity and criticality}

For a proper coloring of $C$, let $q$ be its number of distinct colors
and let $t$ be the number of triples $T_i$ that use exactly two colors.
Each triple uses either two or three colors because it contains the
cycle edge $x_ix_{i+1}$.

\begin{lemma}\label{lem:cycle}
For every proper coloring of $C$,
\begin{equation}\label{eq:cyclebound}
 t\leq6\qquad\text{and}\qquad q+t\leq9.
\end{equation}
\end{lemma}

\begin{proof}
The nonedge pairs in $T_i$ are $\{x_i,x_{i+3}\}$ and
$\{x_{i+1},x_{i+3}\}$.  As $i$ varies, these are exactly the fourteen
nonedges of $C$, each occurring once: the first type has cyclic
distance three and the second distance two.  Since a triple contains
an edge, it is two-colored precisely when one of its nonedge pairs is
monochromatic.  Its other nonedge pair cannot also be monochromatic.
Consequently, if $\mathcal C$ is the set of cycle color classes, then
\[
 t=\sum_{A\in\mathcal C}\binom{|A|}{2}.
\]
Every color class is independent and has at most three vertices.  If
$a_j$ counts classes of size $j$, then
\[
 7=a_1+2a_2+3a_3,\qquad q=a_1+a_2+a_3,\qquad
 t=a_2+3a_3.
\]
Hence $q+t=7+a_3\leq9$, since $a_3\leq2$.  Also $t\leq6$:
when $a_3=0,1,2$, respectively, the above first equality bounds $t$
by $3,5,6$.
\end{proof}

\begin{proposition}\label{prop:critical}
For every $k\geq4$, the graph $H_k$ is fully $k$-critical: it has
chromatic number $k$, and every proper subgraph is $(k-1)$-colorable.
\end{proposition}

\begin{proof}
In a $(k-1)$-coloring of $H_k$, each $Q_i$ would need $k-3$
different colors outside the colors on $T_i$.  Thus all seven triples
would have at most two colors, contradicting $t\leq6$.  Conversely,
three-color the cycle and give each $Q_i$ any $k-3$ colors absent from
$T_i$ out of a $k$-color palette.  This proves $\chi(H_k)=k$.

For criticality, color each edge deletion using $k-1$ colors.  Upon
deleting a cycle edge, two-color the resulting seven-vertex path
alternately.  All triples then use at most two colors, leaving at
least $k-3$ colors for each clique.

For the other edges, the proper cycle coloring
\begin{equation}\label{eq:pattern}
 (\col(x_0),\ldots,\col(x_6))=(0,1,0,2,1,0,1)
\end{equation}
has $T_0$ as its only three-color triple.  Rotate this pattern so that
$T_i$ is the exceptional triple.  Every other clique can use $k-3$
colors absent from its two-color triple.  If the deleted edge joins
$y\in Q_i$ to $x_j\in T_i$, give $y$ the color of $x_j$, and give the
remaining $k-4$ vertices of $Q_i$ distinct colors outside the three
colors of $T_i$.  If the deleted edge is $uv$ inside $Q_i$, give $u,v$
the same color outside $T_i$ and give the other $k-5$ vertices
distinct remaining colors outside $T_i$.  This last case occurs only
when $k\geq5$, so such a color exists.  In each case the stated
coloring is proper after deleting the specified edge.

There are no isolated vertices.  For a vertex $v$, restrict a
$(k-1)$-coloring after deletion of any incident edge to $H_k-v$.
Every proper subgraph is contained in an edge or vertex deletion, so
the claim follows.
\end{proof}

\section{A reduction to six cases}

\begin{lemma}\label{lem:extension}
Suppose an ordered $k$-coloring of $H_k$ has a color $c$ absent from
$C$ and present in every $Q_i$.  Delete the seven vertices of color
$c$ and remove that rank from the palette.  This gives an ordered
$(k-1)$-coloring of $H_{k-1}$.  Every full increasing path of the
smaller coloring extends, injectively, to a full increasing path of
the original coloring.
\end{lemma}

\begin{proof}
Each $Q_i$ is a clique, so it has exactly one vertex of color $c$.
Their deletion reduces all seven clique orders from $k-3$ to $k-4$;
the cycle and every attachment relation remain unchanged.

Every edge of the smaller graph has both ends in a block
$B_i=Q_i\cup T_i$.  For an edge in or incident with a clique this is
immediate.  The cycle edge $x_ix_{i+1}$ belongs to $B_i$.  Further,
every vertex belongs to at least one block.

Take a full increasing path $P$ of the smaller coloring.  If $c$ is
an interior rank, the vertices immediately below and above $c$ on
$P$ are consecutive and therefore adjacent.  Choose a block $B_i$
containing this edge.  Its deleted color-$c$ clique vertex is adjacent
to both ends of the edge; insert it between them.  If $c$ is an
extreme rank, choose a block containing the first or last vertex of
$P$, respectively, and insert its deleted vertex at that end.  The
result is a simple full increasing path in $H_k$.  Fix the least
eligible block index to make the extension a function.  Removing the
unique color-$c$ vertex recovers $P$, so distinct paths have distinct
extensions.
\end{proof}

\begin{lemma}\label{lem:common}
Every ordered $k$-coloring of $H_k$, $k\geq10$, has a color absent from
$C$ and present in every $Q_i$.
\end{lemma}

\begin{proof}
Let $q,t$ be as in Lemma~\ref{lem:cycle}.  There are $k-q$ colors
absent from $C$.  A three-color $T_i$ forces $Q_i$ to use all $k-3$
colors outside that triple.  A two-color $T_i$ allows $k-2$ colors
and its clique omits exactly one of them.  Thus at most $t$ colors
absent from $C$ can be omitted by at least one clique.  By
\eqref{eq:cyclebound}, $k-q>t$ when $k\geq10$, proving the claim.
\end{proof}

Every ordered $h$-coloring of an $h$-chromatic graph has at least one
full increasing path.  Indeed, orient its edges toward larger colors,
and assign to each vertex the number of vertices in a longest directed
path ending there.  These numbers properly color the graph.  If no
path had $h$ vertices, they would give an $(h-1)$-coloring, a
contradiction.  In the color orientation a path has at most $h$
vertices, and a path with $h$ vertices uses every rank.

It follows from Lemmas~\ref{lem:extension} and \ref{lem:common} that
if some coloring of $H_k$, $k\geq10$, had exactly one full increasing
path, then its restricted coloring of $H_{k-1}$ would have at most
one.  By Proposition~\ref{prop:critical} and the preceding paragraph
it would have at least one, hence exactly one.  Repeating this
argument reaches $H_9$.  We now rule out such a coloring for every
$H_k$ with $4\leq k\leq9$.

\section{Finite certificates}\label{sec:certificates}

For reproducibility, label $x_i$ by $i$, $0\leq i<7$, and label the
vertices of $Q_i$ consecutively from $7+i(k-3)$ to
$7+(i+1)(k-3)-1$.  For each $k=4,\ldots,9$, the accompanying CNF asks
whether $H_k$ has an ordered coloring with exactly one full
increasing path.  Its graph edges are generated from
\eqref{eq:triples}.  We describe both the path encoding and every
symmetry restriction used.

Let $C(v,j)$ mean $\col(v)=j$, where $0\leq j<k$.  The formula gives
each vertex exactly one color and makes adjacent vertices differently
colored.  For each vertex $v$ and rank $j$, let $R(v,j)$ assert that
at least one path colored $0,\ldots,j$ ends at $v$, and let $M(v,j)$
assert that at least two distinct such paths end at $v$.  At the first
rank, the exact semantics are
\[
 R(v,0)\longleftrightarrow C(v,0),\qquad \neg M(v,0).
\]
For $j>0$, writing $N(v)$ for the neighbors of $v$, impose
\begin{align}
 R(v,j)&\longleftrightarrow C(v,j)\land
           \bigvee_{u\in N(v)}R(u,j-1),\label{eq:reach}\\
 M(v,j)&\longleftrightarrow C(v,j)\land\left[
    \begin{aligned}
       &\bigvee_{u\in N(v)}M(u,j-1)\\
       {}\lor{}&\bigvee_{\substack{u,w\in N(v)\\u<w}}
                 \bigl(R(u,j-1)\land R(w,j-1)\bigr)
    \end{aligned}\right].
           \label{eq:multiple}
\end{align}
Paths entering $v$ from different neighbors are distinct, and
strictly increasing colors prohibit a repeated vertex.  Therefore
these recurrences exactly count the relevant path prefixes up to
two, including multiple prefixes from one neighbor.  Every gate in
the CNF is defined in both directions.  Finally, the formula
requires exactly one $R(v,k-1)$ and no $M(v,k-1)$.  A sequential
at-most-one counter enforces the first condition.  These constraints
express exactly one full path.

The following restrictions preserve completeness.  Vertices within a
$Q_i$ are true twins, so a clique automorphism sorts their colors in
increasing numeric vertex order.  Rotation of the cycle and attached
cliques permits $x_0$ to be a minimum-color cycle vertex.  Finally
one may require
\begin{equation}\label{eq:lowhalf}
 \col(x_0)\leq\left\lfloor\frac{k-1}{2}\right\rfloor.
\end{equation}
To see this, let $a$ and $b$ be the minimum and maximum colors on the
cycle.  Either $a$ or $k-1-b$ is at most the right side of
\eqref{eq:lowhalf}.  If necessary reverse the entire palette order:
reversal bijects full increasing paths with their reversed paths.
Then rotate a cycle minimum to $x_0$ and sort the cliques again.
Arbitrary permutations of color ranks are not used, since they need
not preserve increasing paths.

\begin{proposition}[Certified finite cases]\label{prop:finite}
Every ordered $k$-coloring of $H_k$ has at least two full increasing
paths for $4\leq k\leq9$.
\end{proposition}

\begin{proof}
The formulas just specified are unsatisfiable.  Table~\ref{tab:checks}
reports their sizes and the number of checked lemma additions in the
positive-hint LRAT refutations.  The supplementary archive contains
the exact CNFs, proof files, and a standard-library Python verifier.
It rebuilds the graph and each CNF independently of the discovery
encoder, compares the clauses in exact order, and checks every LRAT
addition by unit propagation using its listed positive hints.
Negative RAT hints are rejected.  Deletion records can be ignored:
retained clauses are independently proved consequences of the
original formula.  In each of the six cases the checked proof derives
the empty clause.

\begin{table}[ht]
\centering
\caption{Finite formulas and independently checked refutations.}
\label{tab:checks}
\begin{tabular}{rrrrrr}
\toprule
$k$& vertices& edges& variables& clauses& LRAT additions\\
\midrule
4&14&28&622&2,274&360\\
5&21&56&1,791&7,205&1,734\\
6&28&91&4,066&16,972&4,003\\
7&35&133&7,951&33,738&6,913\\
8&42&182&14,034&60,002&11,699\\
9&49&238&22,987&98,599&16,194\\
\bottomrule
\end{tabular}
\end{table}

By the exact encoding and the sound symmetry restrictions, a coloring
with exactly one full increasing path would satisfy the corresponding
CNF, contrary to its refutation.  A full path exists by the argument
before this section, so every coloring has at least two.
\end{proof}

The companion Lean~4.35.0 development also proves unsatisfiability of
the six embedded formulas using the standard-library LRAT soundness
theorem.\footnote{These proofs use \texttt{decide +native} and rely on
Lean's compiler and native-evaluation infrastructure.  Basic adjacency
properties and the Boolean cycle computation are also proved in Lean.}
The encoding and symmetry arguments, criticality, the uniform reduction,
and the orientation--coloring criterion remain outside the formal
development.  Thus the Lean results concern the finite formulas and
do not constitute a formalization of Theorem~\ref{thm:main}.

As a separate finite check, a direct path-counting program enumerates
all 41,664 colorings of $H_4$ up to permutations inside the cliques
and 25,452,840 \emph{irreducible} colorings of $H_5,\ldots,H_9$,
again up to clique permutations.  Here irreducible means that no
cycle-absent color occurs in every $Q_i$; by
Lemma~\ref{lem:extension}, every omitted coloring inherits two
paths from the preceding case.  The program found no coloring with
fewer than two full paths.  This check uses neither the CNF encoder
nor the LRAT traces.

\begin{theorem}\label{thm:main}
For every integer $k\geq4$, the graph $H_k$ is fully $k$-critical and
has no minimal orientation.
\end{theorem}

\begin{proof}
Criticality is Proposition~\ref{prop:critical}.  Proposition
\ref{prop:finite} gives at least two full increasing paths in every
ordered coloring for $4\leq k\leq9$.  Lemmas~\ref{lem:extension} and
\ref{lem:common}, together with the existence of at least one full
path in every ordered coloring, propagate this assertion to all
$k\geq10$.  The orientation--coloring criterion in the introduction
excludes a minimal orientation.
\end{proof}

The result concerns the specific family recorded in
\cite[\S13.8]{JT}; Youngs's earlier work already disproved the
general assertion that all critical graphs have minimal orientations
\cite{Youngs}.  Equation~\eqref{eq:parameters} also gives the upper
bound $f(k)\leq3k-8$ for the degree threshold $f(k)$ defined in
\cite[\S13.8]{JT}.  In particular, the present graphs yield
$f(6)\leq10$ and $f(7)\leq13$.

\section*{Data availability}
The graph data, formulas, refutations, checking and enumeration code,
and Lean sources are available in the accompanying
\href{https://github.com/alexgocardinal/hanson-toft-minimal-orientations}{GitHub repository},
which pins the Lean version and documents the verification scope.

\begin{thebibliography}{9}
\bibitem{Gallai}
T.~Gallai, On directed paths and circuits, in \emph{Theory of Graphs}
(P.~Erd\H{o}s and G.~Katona, eds.), Academic Press, New York, 1968,
pp.~115--118.

\bibitem{JT}
T.~R.~Jensen and B.~Toft, \emph{Graph Coloring Problems}, Wiley,
New York, 1995, \S13.8, pp.~215--216.
\url{https://doi.org/10.1002/9781118032497.ch13}.

\bibitem{Youngs}
D.~A.~Youngs, Minimal orientations of colour critical graphs,
\emph{Combinatorica} \textbf{15} (1995), 289--295.
\url{https://doi.org/10.1007/BF01200761}.
\end{thebibliography}

\end{document}
